\section{数据结构变种、工程应用和实验专题}

这一组专题补齐数据结构的变种、工程应用和实验要求。它不是题号清单，而是把同一类结构放到刷题、C++ 容器和系统源码之间比较，帮助读者理解为什么同一个复杂度在不同场景下表现差异很大。

\subsection{数组、字符串和连续内存的变种}

数组的优势不是语法简单，而是连续内存、随机访问和下标可计算。刷题里常见的前缀和、差分、双指针、二分、原地交换，都依赖下标能在常数时间跳转。工程里数组还带来缓存局部性：顺序扫描通常比链表快很多，即使二者同为线性复杂度。学习数组题时，要把每个下标变量写成区间语义，例如左闭右开窗口、已处理前缀、未知区、已确认后缀。只要区间语义清楚，代码中的加一减一就有根据。

数组变种很多。前缀和把区间求和变成两个端点相减，适合静态区间查询；差分把区间加法变成两个端点标记，适合批量更新后一次还原；二维前缀和把矩形查询变成四个角组合；二维差分适合批量矩形更新；坐标压缩把大值域映射到紧凑下标，适合树状数组、线段树和计数数组。每个变种都要问一个问题：它把哪类重复工作提前保存了，代价是多了什么空间和初始化步骤。

字符串本质上也是字符数组，但它多了编码、复制和匹配成本。LeetCode 基础字符串题常假设小写字母，此时 \code{std::array<int, 26>} 比哈希表更直接；若扩展到 ASCII，可以用 128 或 256 长度数组；若扩展到 Unicode，就不能简单按字节计数。字符串切片也要注意成本，\code{substr} 可能创建新字符串，热循环里更适合用下标区间、\code{string_view} 或 Trie。面试时如果能主动说明字符集假设，答案会更稳。

数组题的实验很直接。第一，给一个窗口题打印 \code{left}、\code{right}、窗口内容和计数表，观察每次移动只改变一个元素。第二，把二维矩阵随机生成后，用暴力矩形求和和二维前缀和对拍，确认四角公式。第三，把 \code{vector} 和 \code{list} 各存一百万个整数，分别顺序求和，比较时间，理解缓存局部性。第四，故意保存 \code{vector} 元素地址后继续 \code{push_back}，观察扩容后旧地址失效。

\subsection{链表、侵入式节点和稳定身份}

链表题的难点不是没有下标，而是节点身份和连接关系分离。数组元素的位置由下标决定，移动元素会改变位置；链表节点的位置由前后指针决定，重连节点不会改变节点自身地址。这个差异让链表适合 LRU、任务队列、空闲块链、内核对象列表和需要稳定迭代器的场景。刷题里的 \code{ListNode} 是最简模型，真实工程里一个业务对象可能带多个链表节点，分别挂进不同索引。

链表有多种变种。单链表空间少，适合反转、合并、快慢指针和环检测；双链表能从已知节点 O(1) 删除，适合 LRU 和双端队列；循环链表适合轮转和哨兵化边界；跳表在多层链表上提供期望对数查找；侵入式链表把链接字段放进业务对象，减少额外分配。每个变种都在回答访问模式：是否需要前驱，是否需要随机查找，是否需要稳定节点，是否会频繁移动到头尾。

链表面试要能讲副作用。反转链表会改变输入结构，回文链表若反转后半段，工程里最好恢复；合并链表可以复用节点，也可以新建节点，二者所有权不同；删除节点时如果平台管理内存，不应该随意释放输入节点；相交链表比较的是地址不是值；环形链表的快慢指针证明来自相对速度，而不是经验。只要这些边界说不清，代码短也不代表理解深。

链表实验建议做三组。第一，把反转链表每一步的 \code{prev}、\code{cur}、\code{next} 打印成边关系，确认没有丢失剩余链。第二，写 LRU 的 \code{get} 后移动到头部，统计哈希查询和链表移动次数。第三，写一个侵入式节点对象，同时挂入按时间排序链表和按 key 哈希桶，体会同一对象进入多个索引时删除顺序有多重要。

\subsection{栈、队列、双端队列和单调结构}

栈的核心是最近未完成结构。括号匹配保存最近未闭合左括号，表达式求值保存上一层上下文，单调栈保存还在等待未来元素解决的下标。不要把“用栈”当结论，必须说清栈顶代表什么，什么事件触发入栈，什么事件触发出栈，出栈时答案是否已经确定。柱状图最大矩形、每日温度、下一个更大元素、最长有效括号虽然都用栈，但弹出的语义完全不同。

队列的核心是按到达顺序处理。BFS 队列保证距离层次，任务队列保证提交顺序，滑动窗口队列保存候选下标。双端队列进一步允许两端进出，适合 0-1 BFS 和单调队列。单调队列不是普通队列加排序，而是在队尾删除被新元素支配的候选，在队首删除过期候选。每个元素最多进出一次，所以线性复杂度来自摊还分析。

单调结构最容易被写成模板。真正要问的是支配关系：一个旧候选为什么被新候选删除后永远不可能成为答案？滑动窗口最大值里，新值更大且更晚，旧小值在未来窗口中既不更大也不更持久；受限子序列和里，新 DP 值更大且下标更晚，也支配旧候选；柱状图中，遇到更矮柱时，被弹出柱子的右边界已经确定。支配关系说不清，单调结构就不可靠。

实验可以把队列状态打印出来。对滑动窗口最大值，记录每轮右端进入、队尾弹出、队首过期和答案读取。对 0-1 BFS，构造零权边和一权边混合的小图，观察零权边入队首如何保证距离顺序。对表达式求值，打印栈中数字和运算符，观察乘除为什么要立即处理，加减为什么可以延迟。

\subsection{哈希、前缀状态和等价类}

哈希表的本质是把对象归入等价类。两数之和把历史值映射到下标，字母异位词把字符串映射到规范签名，前缀和把区间条件映射成历史前缀，账户合并把邮箱映射到账户集合。哈希不是魔法，它只把等值查询变便宜；如果题目需要顺序、前驱、后继或范围最值，哈希表就不是合适结构。

前缀哈希题要先写代数关系。和为 K 的子数组来自 \code{prefix[j] - prefix[i] == k}；连续子数组和来自两个前缀同余；连续数组来自把 0 当成 -1 后找相同前缀；优美子数组来自奇数计数差。关系不同，哈希表保存内容也不同：计数题保存频次，最长题保存最早下标，判定题保存是否出现，最短题可能需要单调队列而不是普通哈希。

哈希 key 的设计决定正确性。异位词可以排序字符串做 key，也可以用字符计数做 key；计数 key 必须带分隔符，否则 \code{1,11} 和 \code{11,1} 可能混淆；结构化 key 可以用 pair 或自定义结构，但要提供哈希函数和相等比较；浮点数不适合直接做哈希 key，工程里要转成可比较的离散表示。面试中讲 key 比讲容器名字更重要。

实验建议做哈希对拍。第一，写和为 K 的子数组暴力版和哈希版，用包含负数、零和重复前缀的随机数组对拍。第二，把先查询后更新改成先更新后查询，观察 target 为 0 时如何出错。第三，给异位词分组设计没有分隔符的计数 key，构造冲突案例。第四，对 \code{unordered_map} 调用和不调用 \code{reserve}，比较大量插入时 rehash 次数。

\subsection{堆、优先队列和动态候选集}

堆解决的是动态候选集的极值查询。排序能给全局顺序，但每次插入删除后重新排序成本高；堆只维护足够的信息让堆顶是最大或最小。前 K 高频、数据流中位数、合并 K 个链表、任务调度、最小加油次数，都是把“下一步选谁”变成堆顶查询。要注意堆不能回答任意元素删除，也不能直接给有序遍历。

双堆是数据流中位数的核心。大顶堆保存较小一半，小顶堆保存较大一半，大小差不超过一，且大顶堆堆顶不大于小顶堆堆顶。两个不变量缺一不可。若要支持滑动窗口删除，就必须加延迟删除表，删除请求先记录，等元素到堆顶时再真正弹出。这和 Linux RCU 的延迟释放不是同一机制，但都体现“逻辑删除”和“物理清理”分离。

堆题要警惕比较器。C++ \code{priority_queue} 默认大顶堆，小顶堆要用 \code{greater} 或自定义比较器。结构体比较器应该明确谁的优先级低，而不是凭直觉写小于号。若堆中保存 pair，默认比较会先比第一维再比第二维；如果业务只看频次或距离，就要写清楚 tie-breaker。比较器一旦不满足严格弱序，结果可能不稳定。

堆实验可以从三个版本比较。第 K 大元素用排序、大小为 K 的小顶堆和快速选择分别实现，比较复杂度和适用场景。合并 K 个有序链表用每轮扫描 K 个头、堆和分治三种方式实现，观察 N 和 K 对性能的影响。任务调度器同时写模拟堆版本和最高频公式版本，解释公式为什么不需要逐步模拟。

\subsection{树、Trie、平衡树和区间树视角}

树题的关键是状态方向。自顶向下问题把父节点约束传给子节点，例如验证 BST 的上下界；自底向上问题让子树向父节点汇报信息，例如平衡二叉树、最大路径和、直径；按层问题用队列维护层边界，例如层序遍历和最小深度。把方向说清楚，递归或迭代只是实现选择。

BST 的价值是有序性。中序遍历得到升序序列，查找可以按大小排除半棵树，前驱后继可以沿树结构寻找。普通二叉树没有这些性质，不能把 BST 的值域剪枝搬过去。红黑树、AVL、Treap、跳表都是为了在动态更新时保持近似平衡；LeetCode 大多不要求实现平衡树，但 \code{map}、\code{set} 和 Linux rbtree 背后都依赖这一类思想。

Trie 把字符串集合按前缀共享。实现 Trie、单词搜索 II、搜索建议系统、异或最大值位 Trie 都是这个思想。数组子节点适合小固定字符集，哈希子节点适合稀疏大字符集；单词结束标记不能省，因为前缀存在不等于单词存在。Trie 的空间可能很大，工程里会考虑压缩 Trie、Radix Tree 或 DAWG 等变体。

树实验要同时做递归和迭代。用显式栈写中序遍历，比较它和递归版本的状态；用队列写层序遍历，打印每层 size；用 Trie 插入一批单词，统计节点数和共享前缀数；用 \code{std::map} 实现日程安排，理解前驱和后继为什么足以判断新区间冲突。

\subsection{图、状态图和搜索维度}

图题第一步永远是建模。网格题的节点可以是格子，也可以是格子加钥匙集合、剩余障碍次数或时间；课程表的节点是课程，边是依赖；单词接龙的节点是单词，边是一位字符变化；转盘锁的节点是四位状态，边是一次拨动。节点维度少了会漏解，节点维度多了会超时。

BFS 适合无权最短路，因为队列按层扩展，第一次到达就是最短。DFS 适合连通块、路径枚举和拓扑状态检测。拓扑排序适合有向依赖，入度为零代表当前可执行。并查集适合动态连通，但不保存路径。最短路算法选择取决于边权：等权 BFS，非负 Dijkstra，负边 Bellman-Ford 或 SPFA 风险评估，边权只有 0 和 1 则 0-1 BFS。

图的性能常败在邻居生成。单词接龙如果每次和所有单词比较，会很慢；通配模式桶把邻居查询降到共享模式。网格 BFS 如果出队才标记访问，可能被同层多个父节点重复入队；入队时标记能控制规模。课程表如果边方向建反，判环可能仍然对，但输出顺序会错。图题必须把邻接表、访问标记和状态编码讲清楚。

实验建议做状态维度对比。打开转盘锁写普通 BFS 和双向 BFS，统计扩展节点数。障碍消除最短路分别用二维 visited 和三维 visited，对比二维为何错误剪枝。课程表分别按两种边方向建图，观察输出顺序差异。最小体力路径写 Dijkstra、二分 BFS、并查集三版，比较它们利用的是哪种单调性。

\subsection{动态规划状态表和空间压缩}

DP 的起点不是公式，而是暴力搜索树。每个搜索节点描述处理到哪里、剩余资源是什么、已经做了什么选择；如果不同路径到达同一个节点后未来选择完全一样，就有重叠子问题。状态定义就是给这种节点命名。定义不稳定时，代码会偷偷依赖历史，边界也会变成补丁。

线性 DP 看最后一步，二维 DP 看两个前缀或网格位置，背包 DP 看物品阶段和容量，区间 DP 看最后合并点，树形 DP 看子树向父节点汇报的信息，状态压缩 DP 看集合中哪些元素已访问。每一类 DP 都要问：状态是否足够，转移是否覆盖全部选择，遍历顺序是否保证依赖已计算，初始化是否对应空状态。

空间压缩是高风险优化。0-1 背包倒序容量，是为了让当前物品只能用一次；完全背包正序容量，是为了允许当前物品重复使用；LCS 一维压缩需要保存左上角旧值；股票状态机需要先保存昨天状态，避免同一天污染。初学时先写清二维或完整状态，确认正确后再压缩。

DP 实验可以非常具体。打印编辑距离表，手动解释每个格子来自插入、删除还是替换。把 0-1 背包容量循环改成正序，用一个物品重复使用的反例验证错误。把零钱兑换 II 的循环顺序调换，观察组合数如何变成排列数。给股票冷冻期状态机打印每天 hold、sold、rest，观察买入来源为什么受限制。

